\chapter{Checagens preliminares}\label{checagens-preliminares-juxe1-executadas}

Foi implementado o script \texttt{checagens\_preliminares.py}, que usa
somente a biblioteca padrão de Python. A saída está em
\texttt{checagens\_preliminares.json}. Ele verifica cinemática; não
implementa a análise estatística proposta.

\section{Escala de frequência}\label{escala-de-frequuxeancia}

O estudo LVK de 2026 reporta
\(m_g\leq1{,}92\times10^{-23}\,\mathrm{eV}/c^2\) a 90\% de credibilidade
em seu teste de dispersão. Seu modelo pressupõe sinal descrito
suficientemente bem por polarizações tensoriais; esse valor não é um
limite universal para toda teoria ou mistura de modos \autocite{lvk2026}.

Usando esse valor apenas como benchmark, obtém-se
\(f_g=4{,}6425\times10^{-9}\,\mathrm{Hz}\):

{
\begingroup
\small
\singlespacing
\setlength{\tabcolsep}{4pt}
\begin{longtable}[]{@{}
  >{\raggedright\arraybackslash}p{(\linewidth - 2\tabcolsep) * \real{0.4286}}
  >{\raggedleft\arraybackslash}p{(\linewidth - 2\tabcolsep) * \real{0.5714}}@{}}
\toprule\noalign{}
\begin{minipage}[b]{\linewidth}\raggedright
Frequência de referência
\end{minipage} & \begin{minipage}[b]{\linewidth}\raggedleft
\(x=(f_g/f)^2\)
\end{minipage} \\ \addlinespace[4pt]
\midrule\noalign{}
\endhead
\bottomrule\noalign{}
\endlastfoot
\(10^{-8}\) Hz, faixa de PTA & \(2{,}1553\times10^{-1}\) \\ \addlinespace[4pt]
\(10^{-7}\) Hz, faixa de PTA & \(2{,}1553\times10^{-3}\) \\ \addlinespace[4pt]
\(10^{-4}\) Hz, referência de baixa frequência para LISA &
\(2{,}1553\times10^{-9}\) \\ \addlinespace[4pt]
\(100\) Hz, interferometria terrestre & \(2{,}1553\times10^{-21}\) \\ \addlinespace[4pt]
\(3000\) Hz, referência do TG para detector ressonante &
\(2{,}3948\times10^{-24}\) \\ \addlinespace[4pt]
\end{longtable}
\endgroup
}

Esses números justificam investigar dispersão local em PTAs. \textbf{Não
são razões universais de amplitudes adicionais nem previsões de
detectabilidade.} Tampouco eliminam testes interferométricos de
dispersão: pequenos efeitos podem acumular fase ao longo da propagação.

Para o benchmark histórico do artigo de 2004,
\(m_gc^2=4{,}4\times10^{-22}\) eV e \(f=1{,}1\times10^{-7}\) Hz,
encontra-se \(x=0{,}93547\), \(v_g/c=0{,}2540\) e
\((c/v_g)^2-1\simeq14{,}50\). Isso mostra que esse exemplo não está em
regime quase nulo. É uma checagem de consistência de aproximação, não
uma nova descoberta sobre polarizações.

\section{Vínculos e matriz de
marés}\label{vuxednculos-e-matriz-de-maruxe9s}

Foram calculados postos de matrizes com aritmética racional exata em
quatro configurações de frequência e número de onda. Nos três casos
massivos, os vínculos linearizados considerados deixam seis amplitudes
no modelo histórico e cinco em Fierz--Pauli, com postos respectivos seis
e cinco no mapa para \(R_{0i0j}\). No caso nulo testado, o posto
observável é dois. Foram verificadas a relação escalar acima, simetrias
de Riemann, Bianchi e anulação da curvatura de uma perturbação de
coordenadas linear.

Isso valida os benchmarks implementados, não demonstra estabilidade,
acoplamento a fontes ou recuperação não linear da RG.

\section{Informação a priori no benchmark de
MeerKAT}\label{informauxe7uxe3o-a-priori-no-benchmark-de-meerkat}

Com \(T=4{,}5\) anos julianos, a escolha \(m_{\max}c^2=h_{\mathrm P}/T\)
corresponde a \(2{,}9123\times10^{-23}\) eV. Uma priori uniforme entre
zero e esse máximo já possui percentil de 90\% igual a
\(2{,}6210\times10^{-23}\) eV, antes de observar dados.

Os limites de 90\% reportados por Zhao e Wang são \(2{,}10\), \(2{,}58\)
e \(2{,}25\times10^{-23}\,\mathrm{eV}/c^2\) nas três configurações. A
comparação aritmética, particularmente para a configuração ER, motiva
medir quanto a posterior realmente acrescenta à priori. \textbf{Essa
proximidade, por si só, não prova que a inferência seja inválida ou que
os dados não contenham informação.} Serão necessárias distribuições
completas e testes controlados \autocite{zhao2026}.

A execução das checagens é descrita no apêndice de reprodutibilidade e no README do
repositório.
